The proposed detector is positioned qualitatively rather than as numerically superior to earlier blink detectors because direct numerical ranking would be misleading given differences in datasets, blink definitions, sampling rates, and scoring procedures, as summarised in Table~\ref{tab:literature_comparison}. In particular, matching criteria ranged from sample-level classification to event-level detection, so apparently similar performance measures did not necessarily represent the same detection requirement. This heterogeneity applies to threshold-based and unsupervised approaches alike~\citep{chang2016detection,tran2021detection,kleifges2017blinker,agarwal2019blink,jefri2022eye,cao2021unsupervised,rashida2023quantitative}. Its performance was competitive with interpretable single-channel threshold detectors while being assessed under a strict event-level overlap criterion on both independent driving corpora. This evaluation required each predicted region to correspond adequately to a reference blink event rather than rewarding isolated correctly labelled samples. The detector also avoided the dataset-specific threshold tuning commonly used by comparable approaches, which indicates that its relative performance was obtained under a more transferable operating procedure rather than through adaptation to a particular recording set.
